\documentclass[11pt]{article}

\usepackage[utf8]{inputenc}
\usepackage[T1]{fontenc}
\usepackage{lmodern}
\usepackage{geometry}
\usepackage{amsmath,amssymb,amsfonts,amsthm,mathtools}
\usepackage{bm}
\usepackage{enumitem}
\usepackage{microtype}
\usepackage{hyperref}
\usepackage{titlesec}
\usepackage{booktabs}
\usepackage{array}
\usepackage{setspace}
\usepackage{mathrsfs}
\usepackage{physics}

\geometry{margin=1in}
\hypersetup{colorlinks=true, linkcolor=black, urlcolor=blue, citecolor=black}
\setlist[enumerate]{topsep=0.4em,itemsep=0.35em}
\setlist[itemize]{topsep=0.3em,itemsep=0.25em}
\titleformat{\section}{\large\bfseries}{\thesection.}{0.5em}{}
\titleformat{\subsection}{\normalsize\bfseries}{\thesubsection.}{0.5em}{}
\setstretch{1.06}

\newcommand{\Aether}{\AE{}ther}
\newcommand{\Aflow}{\AE{}ther-flow}
\newcommand{\TheoryName}{\Aflow{} Interpretation of Relativity}
\newcommand{\TheoryProgram}{\Aether{} / \Aflow{} framework}
\newcommand{\CompletedMetric}{g^{\sharp}}
\newcommand{\PreProtoPackage}{\mathfrak Q^{\mathrm{preproto}}}
\newcommand{\PreProtoCore}{\mathfrak Q^{\mathrm{core}}}

\newtheorem{definition}{Definition}
\newtheorem{proposition}{Proposition}
\newtheorem{remark}{Remark}
\newtheorem{corollary}{Corollary}
\newtheorem{theorem}{Theorem}

\title{\textbf{The \TheoryName{}}\\[0.4em]
\large Primitive-Reservoir Observer-Reduction\\
\large Proto-Germ-Generating Substrate Pre-Proto-Germ\\
\large Observer-Relevant Core Next Benchmark Criterion}
\author{}
\date{}

\begin{document}

\maketitle

\begin{abstract}
The active primitive-reservoir observer-reduction line has now moved below the full pre-proto-germ package itself. The pre-proto-germ dynamics theorem already removed the proto-germ package as the primitive stopping point, and the later observer-relevant core collapse theorem then spent the next honest internal compression move available there: within the full pre-proto-germ package \(\PreProtoPackage\), the live benchmark-facing burden collapses to the smaller observer-relevant pre-proto-germ core \(\PreProtoCore\).

This note records the routing consequence of that theorem. Another same-output deeper-origin relay that leaves \(\PreProtoCore\) unchanged no longer counts as the honest default next benchmark-facing manuscript, and neither does a reopening of older screened layers. The next honest benchmark-facing move must therefore either derive or justify \(\PreProtoCore\) from still more primitive substrate structure in a way that changes benchmark-facing status, or prove a still sharper theorem-level collapse strictly inside \(\PreProtoCore\).

The benchmark boundary remains fixed throughout. The one operative metric is still exactly \(\CompletedMetric\), the ordinary matter route is unchanged, there is no second operative metric, the low-energy matter law is not enlarged, and no stronger promotion language is licensed. The positive result of this note is therefore routing discipline at the compressed pre-proto-germ core frontier.
\end{abstract}

\section{Introduction}

The active derivational program of \TheoryName{} is no longer at a stage where the full pre-proto-germ package should be treated as the default unexplained stopping point. The theorem deriving the current proto-germ package from a more primitive pre-proto-germ package already moved the frontier below the proto-germ layer. The later core-collapse theorem then sharpened that frontier again by proving that, inside the full pre-proto-germ package, only a smaller observer-relevant core remains benchmark-facingly live.

The present note records the routing consequence of that sharpening. It does not reopen the pre-proto-germ mathematics, and it does not claim a completed first-principles substrate derivation of gravity. It records only which kind of next manuscript would now count as an honest benchmark-facing gain below the compressed pre-proto-germ core frontier.

\paragraph{Framework and claim boundary.}
\TheoryProgram{} denotes the broader \Aether{} / \Aflow{} framework built on the claims that the \Aether{} is the underlying four-dimensional substrate of reality, the \Aflow{} is its intrinsic ordered motion, observed three-dimensional space is the local experiential slice of that deeper substrate, S-time is the experienced order of change arising from matter, light, and the \Aflow{}, and observed expansion is the three-dimensional appearance of deeper four-dimensional ordered motion. Within that framework, \TheoryName{} denotes the exact-closure theory in which the effective gravitational dynamics are adopted to be exactly Einsteinian with universal matter coupling. In the active sequence, \emph{adoption} means use of that established relativistic dynamics without claiming substrate derivation, while \emph{derivation} is reserved for a first-principles recovery from explicit substrate variables.

The present note stays strictly inside that boundary. It records only which kind of next manuscript would now count as an honest benchmark-facing gain at the compressed pre-proto-germ core frontier.

\section{Fixed Inputs}

\begin{definition}[Compressed pre-proto-germ core frontier inputs]
The criterion recorded here uses the following fixed inputs.
\begin{enumerate}
    \item The benchmark exact-closure package, which fixes the public one-metric GR target of the repository.
    \item The benchmark gatekeeping layer, including the recorded safeguard that blocks same-output deeper-origin relays from counting as new front-line progress once the live observer-relevant core is unchanged.
    \item The theorem deriving the current proto-germ package from a more primitive pre-proto-germ package \(\PreProtoPackage\).
    \item The later observer-relevant pre-proto-germ core collapse theorem, which compresses the live benchmark-facing burden inside \(\PreProtoPackage\) to the smaller core \(\PreProtoCore\).
\end{enumerate}
\end{definition}

\begin{remark}
The present note does not replace those manuscripts. It records the routing consequence of reading them together under the active safeguard against recursive depth drift.
\end{remark}

\section{Why the Core-Collapse Theorem Is the Frontier}

\begin{definition}[Benchmark-neutral deeper relay at core scope]
A \emph{benchmark-neutral deeper relay at core scope} is a manuscript that preserves the same benchmark one-metric output, the same ordinary matter route, and the same observer-relevant pre-proto-germ core \(\PreProtoCore\) while merely relocating the unexplained substrate layer deeper, restating the full pre-proto-germ package \(\PreProtoPackage\), or repackaging the omitted lower completion data without providing a new compression, necessity, uniqueness, or comparably sharp routing gain.
\end{definition}

\begin{proposition}[The core-collapse theorem is the live frontier]
At the current stage of the primitive-reservoir observer-reduction line, the observer-relevant pre-proto-germ core collapse theorem remains the live benchmark-facing frontier.
\end{proposition}

\begin{proof}
That theorem changes benchmark-facing status at current scope because it proves that the full pre-proto-germ package is not the remaining live burden. The benchmark-facing task now factors through the smaller observer-relevant core \(\PreProtoCore\). By contrast, a later manuscript that preserves the same benchmark package and the same core while merely moving the unexplained layer deeper does not alter benchmark-facing status unless it adds a comparably sharp gain. Therefore the live frontier is now the core-collapse theorem rather than the earlier full-package derivation theorem or any later same-output relay that leaves \(\PreProtoCore\) unchanged.
\end{proof}

\begin{corollary}[Status of the prior pre-proto-germ dynamics theorem]
The theorem deriving the current proto-germ package from the full pre-proto-germ package remains preserved benchmark-facing main-line history, but under the current routing safeguard it is no longer by itself the default current frontier.
\end{corollary}

\begin{proof}
That theorem matters because it removed the proto-germ package as the unexplained stopping point and moved the line into pre-proto-germ scope. But the later core-collapse theorem shows that even the full pre-proto-germ package is no longer the present live benchmark-facing burden. The earlier theorem is preserved history and handoff, while the later theorem defines the current frontier.
\end{proof}

\section{The Next Benchmark Criterion}

\begin{definition}[Admissible next benchmark-facing manuscript at compressed core scope]
At the compressed pre-proto-germ core frontier, a manuscript counts as an admissible next benchmark-facing move only if it does at least one of the following:
\begin{enumerate}
    \item derives or justifies the observer-relevant pre-proto-germ core \(\PreProtoCore\) from still more primitive substrate structure in a way that does more than merely relocate the unexplained layer deeper;
    \item proves a still sharper theorem-level collapse strictly inside \(\PreProtoCore\);
    \item does either of the previous two while preserving the same benchmark one-metric package, the same ordinary matter route, and the same claim boundary.
\end{enumerate}
\end{definition}

\begin{theorem}[Next honest benchmark-facing task at compressed core scope]
The next honest benchmark-facing manuscript below the current compressed pre-proto-germ core frontier must either derive or justify the observer-relevant pre-proto-germ core from still more primitive substrate structure in a benchmark-relevant way, or else prove a still sharper collapse inside that core. A manuscript that only repeats the same core through another deeper relay, or only re-expands the discussion back to the full pre-proto-germ package without such a new gain, does not satisfy the criterion.
\end{theorem}

\begin{proof}
By Proposition 1, the live frontier already lies at the observer-relevant core rather than at the full pre-proto-germ package. Therefore a subsequent manuscript must improve on that status rather than merely accompany it. A deeper-origin manuscript can do so only if it changes benchmark-facing status by more than depth alone. If it simply reproduces the same core through another relay, the routing rule classifies it as benchmark neutral. If it merely restates the wider package whose extra freedom has already been shown benchmark-neutral, it likewise does not change current status. The remaining honest alternatives are therefore exactly those stated in the definition: either a more primitive derivation or justification of the core itself, or a sharper internal collapse of that core.
\end{proof}

\section{What the Next Move Should Not Be}

\begin{proposition}[Screened next moves]
At the compressed pre-proto-germ core frontier, none of the following is the default next benchmark-facing move:
\begin{enumerate}
    \item another same-output deeper-origin relay that leaves \(\PreProtoCore\) unchanged;
    \item another full-package pre-proto-germ restatement or lower-completion relay that reintroduces benchmark-neutral data already screened out by the core-collapse theorem;
    \item a reopening of older screened shell, lift, support, reservoir, ground, stationary-reduction, stationary-Hessian, spectral, or selector layers;
    \item a manuscript that introduces a second operative metric, enlarges the ordinary matter law, or uses stronger promotion language.
\end{enumerate}
\end{proposition}

\begin{proof}
Items (1) and (2) fail the preceding criterion because they do not directly address the current core-level burden. They revisit a same-output relay or a previously compressed larger package rather than removing the remaining freedom in \(\PreProtoCore\) or proving a smaller internal core. Item (3) likewise does not directly address the live burden at current scope. It reopens older screened layers as if they were still frontier-defining rather than settling the present core-level task. Item (4) violates the fixed claim boundary of the benchmark package and therefore cannot count as a disciplined next step on the active line. The benchmark package remains the exact one-metric GR package already fixed by adoption, so any admissible next manuscript must preserve that boundary.
\end{proof}

\begin{remark}
This screening statement is not a denial that those deeper or older manuscripts exist or matter historically. It is a routing statement: they are not the default way to spend the next benchmark-facing move from the compressed pre-proto-germ core frontier.
\end{remark}

\section{Conclusion}

The positive result of this note is routing discipline at the compressed pre-proto-germ core frontier. The observer-relevant pre-proto-germ core collapse theorem remains the live benchmark-facing gain because it already spends the honest internal compression move available below the full pre-proto-germ package: the live burden now lies at the smaller core \(\PreProtoCore\), not at the full package that contains it.

The scope remains narrow and honest. The benchmark package is unchanged: the one operative metric is still exactly \(\CompletedMetric\), the ordinary matter route is unchanged, there is no second operative metric, and the low-energy matter law is not enlarged. Nothing here upgrades adoption to derivation or licenses stronger promotion language.

The next open burden is therefore clear. The next benchmark-facing manuscript should derive or justify the observer-relevant pre-proto-germ core from still more primitive substrate structure in a way that changes benchmark-facing status, unless the sharper honest gain available instead is a still smaller theorem-level collapse strictly inside that core. Another same-output deeper relay that leaves that core unchanged, or a reopening of older screened layers, is not the clean next manuscript target at current scope.

\end{document}
